%========================
% main.tex
%========================
\documentclass[12pt,a4paper,oneside]{report}

% Encoding and language
\usepackage[utf8]{inputenc}
\usepackage[T1]{fontenc}
\usepackage[brazil,english]{babel}
\usepackage{csquotes}

% Typography
\usepackage{mathptmx}
\usepackage{setspace}
\onehalfspacing

% Math and theorem environments
\usepackage{amsmath,amssymb,amsfonts,mathtools,bm}
\usepackage{amsthm}
\newtheorem{theorem}{Theorem}[chapter]
\newtheorem{definition}{Definition}[chapter]

% Figures and tables
\usepackage{graphicx}
\usepackage{float}
\usepackage{placeins}
\usepackage{subcaption}
\usepackage{booktabs}
\usepackage{multirow}
\usepackage{longtable}
\usepackage{array}
\usepackage{tabularx}
\usepackage{ragged2e}
\newcolumntype{Y}{>{\RaggedRight\arraybackslash}X}
\graphicspath{{figures/}}

% Algorithms
\usepackage{algorithm}
\usepackage{algpseudocode}

% Utilities
\usepackage{xcolor}
\usepackage{siunitx}
\usepackage{geometry}
\geometry{a4paper, margin=3cm}

% Hyperlinks
\usepackage[hidelinks]{hyperref}
\usepackage[nameinlink,noabbrev]{cleveref}

% TOC
\usepackage{tocloft}
\renewcommand{\contentsname}{TABLE OF CONTENTS}
\renewcommand{\listfigurename}{LIST OF FIGURES}
\renewcommand{\listtablename}{LIST OF TABLES}

% Chapter formatting
\setcounter{secnumdepth}{2}
\setcounter{tocdepth}{2}

% Custom commands
\newcommand{\thesistitle}{Learning Node Orderings for Autoregressive Graph Generation}
\newcommand{\thesisauthor}{Alysson Amaral da Silva}
\newcommand{\thesisyear}{2026}
\newcommand{\thesiscity}{Curitiba}
\newcommand{\advisorname}{Prof.\ Dr.\ Andr\'e Lu\'is Vignatti}
\newcommand{\R}{\mathbb{R}}
\newcommand{\E}{\mathbb{E}}
\newcommand{\NLL}{\mathcal{L}_{\mathrm{NLL}}}
\newcommand{\policy}{\pi_{\theta}}
\newcommand{\generator}{p_{\phi}}
\newcommand{\OrderSet}{\mathcal{O}}
\newcommand{\GraphSet}{\mathcal{D}}
\newcommand{\maybeincludegraphics}[2][]{%
\IfFileExists{#2}{\includegraphics[#1]{#2}}{\fbox{\parbox{0.88\linewidth}{\centering Missing figure file: \texttt{\detokenize{#2}}}}}%
}

\begin{document}
\selectlanguage{english}

%========================================
% Title page
%========================================
\begin{titlepage}
    \thispagestyle{empty}
    \centering

    {\Large \thesisauthor \par}

    \vspace{3cm}

    {\bfseries\LARGE \MakeUppercase{\thesistitle}\par}

    \vspace{2cm}

    {\large
    Master's Dissertation submitted to the Graduate Program in Informatics,
    Department of Exact Sciences, Federal University of Paran\'a,
    as a partial requirement for the degree of Master in Informatics.\par}

    \vspace{1cm}

    {\large Advisor: \advisorname\par}

    \vfill

    {\large \thesiscity\par}
    {\large \thesisyear\par}
\end{titlepage}

%========================================
% Front matter
%========================================
\pagenumbering{roman}
\setcounter{page}{1}

\begin{abstract}
Sequential graph generative models transform graph-structured objects into ordered construction sequences before learning. This transformation is necessary in autoregressive generators, but it is also a modeling choice: the same graph can induce many different action sequences, and each sequence changes the conditioning context seen by the model. This dissertation studies node ordering as a central component of autoregressive graph generation rather than as a neutral preprocessing step. The work first revisits a autoregressive graph model, reproducing its graph-to-sequence mechanism and using it as a controlled environment for analyzing how deterministic node orderings affect structural reconstruction. The experiments compare degree ordering, breadth-first search, depth-first search, degeneracy, LexBFS, maximum cardinality search, spectral ordering, and random permutations across cycles, trees, binary trees, stars, and Barab\'asi--Albert graphs.

The main contribution is a learned-ordering formulation in which node sequencing is modeled as a graph-conditioned Markov Decision Process. A policy parameterized by a Graph Neural Network selects nodes one at a time, producing an ordering that is translated into an action sequence for an autoregressive model. Because the ordering is discrete and its effect is only observable through downstream generation quality, the policy is optimized with a policy-gradient objective inside an alternating procedure that updates both the ordering controller and the graph generator. The final version of the method, called Attention RL ordering, strengthens the controller with an attention-based graph encoder, three graph layers, and a degree-based warm start. On Barab\'asi--Albert graphs, the method achieves perfect valid-size ratio, near-perfect connectivity, realistic graph size, and competitive hub-oriented statistics, showing that learned orderings can become strong alternatives to handcrafted serializations. Overall, the dissertation argues that node ordering should be treated as an optimization target in sequential graph generation and that reinforcement learning provides a natural framework for coupling this target to downstream generative performance.

\textbf{Keywords:} graph generation; autoregressive models; node ordering; reinforcement learning; graph neural networks.
\end{abstract}

\clearpage

\begin{otherlanguage*}{brazil}
\renewcommand{\abstractname}{Resumo}
\begin{abstract}
Modelos generativos sequenciais de grafos transformam objetos estruturados em grafos em sequ\^encias ordenadas de constru\c{c}\~ao antes do aprendizado. Essa transforma\c{c}\~ao \`e necess\'aria em geradores autoregressivos, mas tamb\'em \'e uma escolha de modelagem: o mesmo grafo pode induzir muitas sequ\^encias de a\c{c}\~oes diferentes, e cada sequ\^encia altera o contexto condicional observado pelo modelo. Esta disserta\c{c}\~ao estuda a ordena\c{c}\~ao de n\'os como um componente central da gera\c{c}\~ao autorregressiva de grafos, e n\~ao como uma simples etapa neutra de pr\'e-processamento. O trabalho revisita inicialmente um modelo de grafos autoregressivo, reproduzindo seu mecanismo de convers\~ao de grafos em sequ\^encias e utilizando-o como ambiente controlado para analisar como diferentes ordena\c{c}\~oes determin\'isticas afetam a reconstru\c{c}\~ao estrutural. Os experimentos comparam ordena\c{c}\~ao por grau, busca em largura, busca em profundidade, degeneresc\^encia, LexBFS, \emph{maximum cardinality search}, ordena\c{c}\~ao espectral e permuta\c{c}\~oes aleat\'orias em ciclos, \''arvores, \''arvores bin\'arias, estrelas e grafos de Barab\'asi--Albert.

A principal contribui\c{c}\~ao \`e uma formula\c{c}\~ao aprendida de ordena\c{c}\~ao, na qual o sequenciamento de n\'os \`e modelado como um Processo de Decis\~ao de Markov condicionado ao grafo. Uma pol\'itica parametrizada por uma Rede Neural em Grafos seleciona os n\'os um a um, produzindo uma ordena\c{c}\~ao que \`e traduzida em uma sequ\^encia de a\c{c}\~oes para o modelo autorregressivo. Como a ordena\c{c}\~ao \`e discreta e seu efeito s\'o \`e observado por meio da qualidade gerativa final, a pol\'itica \`e otimizada por gradiente de pol\'itica dentro de um procedimento alternado que atualiza tanto o controlador de ordena\c{c}\~ao quanto o gerador de grafos. A vers\~ao final do m\'etodo, chamada Attention RL ordering, fortalece o controlador com um codificador baseado em aten\c{c}\~ao, tr\^es camadas de grafo e uma inicializa\c{c}\~ao baseada na ordena\c{c}\~ao por grau. Em grafos de Barab\'asi--Albert, o m\'etodo alcan\c{c}a raz\~ao perfeita de tamanho v\'alido, conectividade quase perfeita, tamanho realista e estat\'isticas competitivas relacionadas a hubs, mostrando que ordena\c{c}\~oes aprendidas podem se tornar alternativas fortes \`as serializa\c{c}\~oes constru\'idas manualmente. De forma geral, a disserta\c{c}\~ao defende que a ordena\c{c}\~ao de n\'os deve ser tratada como um alvo de otimiza\c{c}\~ao na gera\c{c}\~ao sequencial de grafos e que o aprendizado por refor\c{c}o oferece uma estrutura natural para acoplar esse alvo ao desempenho gerativo final.

\textbf{Palavras-chave:} gera\c{c}\~ao de grafos; modelos autoregressivos; ordena\c{c}\~ao de n\'os; aprendizado por refor\c{c}o; redes neurais em grafos.
\end{abstract}
\end{otherlanguage*}

\clearpage
\tableofcontents
\clearpage
\listoffigures
\clearpage
\listoftables
\clearpage

\chapter*{LIST OF ACRONYMS}
\addcontentsline{toc}{chapter}{LIST OF ACRONYMS}
\begin{longtable}{p{3cm}p{11cm}}
BA   & Barab\'asi--Albert \\
BFS  & Breadth-First Search \\
DFS  & Depth-First Search \\
DGMG & Deep Generative Model of Graphs \\
GAT  & Graph Attention Network \\
GCN  & Graph Convolutional Network \\
GGNN & Gated Graph Neural Network \\
GNN  & Graph Neural Network \\
GRU  & Gated Recurrent Unit \\
LexBFS & Lexicographic Breadth-First Search \\
MCS  & Maximum Cardinality Search \\
MDP  & Markov Decision Process \\
MLP  & Multi-Layer Perceptron \\
NLL  & Negative Log-Likelihood \\
RL   & Reinforcement Learning \\
\end{longtable}

\clearpage

\chapter*{LIST OF SYMBOLS}
\addcontentsline{toc}{chapter}{LIST OF SYMBOLS}

\noindent\textbf{Graph notation}
\begin{longtable}{p{3.8cm}p{10.2cm}}
$G=(V,E)$              & Undirected, unweighted graph with node set $V$ and edge set $E$ \\
$V$                    & Finite set of nodes \\
$E \subseteq V\times V$& Set of undirected edges \\
$n=|V|$                & Number of nodes \\
$m=|E|$                & Number of edges \\
$A\in\{0,1\}^{n\times n}$ & Adjacency matrix of $G$ (defined after fixing a node ordering) \\
$\mathcal{N}(v)$       & Neighborhood of node $v$ \\
$e_{uv}$               & Edge feature between nodes $u$ and $v$ \\
$G_t$                  & Partial graph at autoregressive construction step $t$ \\
\end{longtable}

\noindent\textbf{Node orderings and sequences}
\begin{longtable}{p{3.8cm}p{10.2cm}}
$\sigma=(v_1,\ldots,v_n)$  & Node ordering; permutation of $V$ \\
$\mathcal{O}_G$            & Set of all admissible node orderings for graph $G$ \\
$\mathcal{T}(G,\sigma)$    & Graph-to-sequence translator; maps $(G,\sigma)$ to an action sequence \\
$a=(a_1,\ldots,a_L)$       & Autoregressive action sequence of length $L$ \\
$a_t$                      & Action taken at step $t$ \\
$a_{<t}=(a_1,\ldots,a_{t-1})$ & Prefix of the action sequence up to step $t$ \\
$L$                        & Total length of an action sequence \\
\end{longtable}

\noindent\textbf{Generator}
\begin{longtable}{p{3.8cm}p{10.2cm}}
$\phi$                       & Parameters of the graph generator \\
$p_\phi$                     & Generator distribution parameterized by $\phi$ \\
$\phi^\star(\theta)$         & Optimal generator parameters for a given policy $\pi_\theta$ \\
$\mathcal{L}_{\mathrm{NLL}}(\phi;G,\sigma)$ & Negative log-likelihood training loss \\
$\mathcal{L}_{\mathrm{val}}(\phi,\mathcal{T}(G,\sigma))$ & Validation negative log-likelihood \\
$\overline{\mathcal{L}}_{\mathrm{val}}(\tau)$ & Average validation loss over $\mathcal{D}_{\mathrm{val}}$ \\
\end{longtable}

\noindent\textbf{Datasets}
\begin{longtable}{p{3.8cm}p{10.2cm}}
$\mathcal{D}$              & Dataset of training graphs \\
$\mathcal{D}_{\mathrm{train}}$ & Training split \\
$\mathcal{D}_{\mathrm{val}}$   & Validation split \\
$N\in[20,60]$              & Target node-count range for Barab\'asi--Albert graphs \\
$m_{\mathrm{BA}}$          & Attachment parameter of the BA preferential-attachment mechanism \\
\end{longtable}

\noindent\textbf{Markov Decision Process}
\begin{longtable}{p{3.8cm}p{10.2cm}}
$\mathcal{S}$              & State space \\
$\mathcal{A}$              & Action space; $\mathcal{A}(s_t)=V\setminus S_t$ denotes available actions at state $s_t$ \\
$\mathcal{P}$              & State transition probability function \\
$\mathcal{R}$              & Step-wise reward function \\
$\gamma$                   & Discount factor (set to $1$ in the episodic setting) \\
$\tau$                     & Complete trajectory (full episode) \\
$s_t=(G,S_t)$              & State at step $t$: input graph $G$ and set of already selected nodes $S_t$ \\
$S_t\subseteq V$           & Set of nodes already selected before step $t$ \\
$u_t$                      & Node selected at ordering step $t$ \\
\end{longtable}

\noindent\textbf{Ordering policy}
\begin{longtable}{p{3.8cm}p{10.2cm}}
$\theta$                   & Parameters of the ordering policy \\
$\theta^\star$             & Optimal policy parameters \\
$\pi_\theta$               & Stochastic ordering policy parameterized by $\theta$ \\
$J(\theta)$                & Expected return of the policy; training objective \\
$J_G(\theta)$              & Expected return of the policy for a single graph $G$ \\
$b(G)$                     & Baseline for variance reduction in REINFORCE \\
$\beta$                    & Entropy regularization coefficient \\
$\mathcal{H}(\cdot)$       & Entropy of a probability distribution \\
\end{longtable}

\noindent\textbf{Rewards}
\begin{longtable}{p{3.8cm}p{10.2cm}}
$R(\tau)$                  & Total delayed reward (trajectory return) \\
$R(G,\sigma,\phi)$         & Reward for ordering $\sigma$ of graph $G$ under generator parameters $\phi$ \\
$R_{\mathrm{pred}}(\tau)$  & Predictive component of the composite reward (based on validation NLL) \\
$R_{\mathrm{struct}}(\tau)$& Structural component of the composite reward (based on generated-sample quality) \\
$\lambda$                  & Trade-off weight between predictive and structural reward \\
\end{longtable}

\noindent\textbf{Graph Neural Network}
\begin{longtable}{p{3.8cm}p{10.2cm}}
$x_v$                      & Input feature vector of node $v$ \\
$h_v^{(k)}$                & Embedding of node $v$ at GNN layer $k$ \\
$h_v^{(T_\sigma)}$         & Final node embedding produced by the policy GNN after $T_\sigma$ layers \\
$h_G^{(T_\sigma)}$         & Graph-level context vector, obtained by pooling node embeddings \\
$\eta^{(k)}$               & Message function at GNN layer $k$ \\
$\psi^{(k)}$               & Node update function at GNN layer $k$ \\
$\square$                  & Permutation-invariant aggregation operator (e.g., sum or mean) \\
$T_\sigma$                 & Number of message-passing rounds in the policy network \\
$z_v^{(t)}$                & Logit score for candidate node $v$ at ordering step $t$ \\
$f_\theta$                 & Scoring MLP that maps node and graph embeddings to a logit \\
$\mathbb{I}\{v\in S_t\}$   & Indicator: equals $1$ if node $v$ has already been selected, $0$ otherwise \\
\end{longtable}

\noindent\textbf{Hyperparameters}
\begin{longtable}{p{3.8cm}p{10.2cm}}
$H$                        & Hidden dimension of the generator \\
$T$                        & Number of DGMG message-passing rounds \\
$\alpha$                   & Learning rate \\
$I$                        & Number of outer iterations in the alternating optimization \\
$K$                        & Number of policy-gradient steps per outer iteration \\
\end{longtable}

\clearpage

%========================================
% Main matter
%========================================
\pagenumbering{arabic}

%========================
% chapters/01_introduction.tex
%========================
\chapter{Introduction}
\label{chap:introduction}

Graphs are a natural representation for systems whose meaning is defined not only by individual entities, but also by the pattern of relations among them. Molecules, social networks, communication systems, power grids, circuits, knowledge graphs, and biological structures are all examples in which topology is part of the object being modeled. This makes graph generation a central task in modern machine learning: a generator should not merely reproduce attributes attached to vertices or edges, but should also learn how local decisions combine into global connectivity, degree distribution, motifs, and structural validity.

Generative models for graphs face a difficulty that is less visible in image or text generation. A graph is invariant to node relabeling, whereas a neural model usually receives a tensor, a sequence, or another ordered representation. When an unordered graph is converted into an ordered training object, the conversion is not unique. The same graph with $n$ nodes admits up to $n!$ node permutations, and each permutation can induce a different sequence of construction actions. Consequently, a graph generator that appears to model a graph distribution may in fact be modeling a distribution over graph serializations. This tension is especially important in autoregressive graph generation, where the model builds a graph step by step and every decision is conditioned on the partial graph already produced.

This dissertation studies that tension through the lens of node ordering. The central argument is that node ordering is not a harmless implementation detail. It determines which vertices are visible when an edge must be predicted, which candidate destinations compete in the model's output distribution, how long-range dependencies appear in the action sequence, and how regular the stop decisions become. In a model such as the Deep Generative Model of Graphs (DGMG) \cite{li2018learning}, the generator alternates among decisions to add a node, add an edge, choose an edge destination, and stop. The ordering of the original graph determines the supervised action sequence used during training. Two orderings of the same graph therefore expose the same topology through different conditional histories, and these histories can vary substantially in difficulty.

The work begins with a controlled reproduction of the graph-to-sequence mechanism and uses this reproduction to compare deterministic orderings across several graph families. Degree ordering, breadth-first search, depth-first search, degeneracy ordering, LexBFS, maximum cardinality search, spectral ordering, and random permutations are evaluated on cycles, trees, binary trees, stars, and Barab\'asi--Albert graphs. This stage shows that different graph families favor different orderings. Cycles benefit from orderings that nearly linearize the ring; trees and binary trees favor level-wise parent-child locality; stars benefit from exposing the hub early; and Barab\'asi--Albert graphs reveal trade-offs among size control, connectivity, and hub structure. These results motivate a stronger claim: because no single handcrafted ordering is uniformly best, ordering should itself become a learnable object.

The main methodological contribution is therefore a deep reinforcement-learning formulation for node ordering. Instead of imposing a fixed ordering rule, a policy observes the input graph and selects one node at a time until a complete permutation is produced. The resulting ordering is translated into a graph generation action sequence, and the quality of this sequence is assessed through downstream generator performance. Since the choice of a node is discrete and the reward is only available after the generator has been trained and evaluated, direct differentiation through the ordering is not possible. Reinforcement learning provides a natural solution: the ordering policy is treated as a stochastic graph-conditioned policy and optimized with a policy-gradient estimator.

The learned-ordering method used in the final experiments is called Attention RL ordering. Its policy network encodes the graph with message passing and attention, maintains the set of already selected nodes, masks invalid actions, and produces a probability distribution over the remaining nodes. The policy is trained in an alternating loop with the autoregressive generator, which in our case of study is the DGMG \cite{li2018learning}, but it could be any other generator; The choice of using DGMG was merely because it is one of the foundational approaches to graph generation through decision sequences. At each outer iteration, the policy samples orderings for the training graphs, the translator produces action sequences, the generator is updated on those sequences, and validation and sampling metrics are used to update the policy. The final version strengthens the earlier controller with three graph layers, an attention mechanism, and a degree-based warm start.

The empirical focus of the final learned-ordering study is the Barab\'asi--Albert family \cite{barabasi1999emergence}. This family is a useful test bed because it combines growth, preferential attachment, hub dominance, sparse connectivity, and variable graph size. Unlike cycles or trees, it does not admit a single exact validity predicate that fully captures structural quality. The evaluation therefore uses several complementary metrics: average generated size, valid-size ratio, connected ratio, degree Gini coefficient, hubiness, and average number of edges. The final Attention RL ordering run achieves perfect valid-size ratio, near-perfect connectivity, and realistic size, while keeping degree inequality and hub-related statistics competitive with the deterministic baselines. This supports the view that a learned ordering policy can act as a flexible structural prior for sequential graph generation.

The dissertation is organized as follows. Chapter~\ref{chap:foundations} introduces the graph-theoretic, generative-modeling, and reinforcement-learning concepts used throughout the work. Chapter~\ref{chap:related-work} positions the dissertation with respect to graph generative models, autoregressive graph generation, ordering-aware methods, and reinforcement learning. Chapter~\ref{chap:dgmg-reproduction} describes the DGMG reproduction and the graph-to-sequence translator that makes the ordering problem explicit. Chapter~\ref{chap:deterministic-orderings} analyzes deterministic orderings and explains how they reshape the action-sequence learning problem. Chapter~\ref{chap:rl-ordering} presents the learned-ordering formulation, including the Markov Decision Process, the policy architecture, the reward, the policy-gradient theorem, and the alternating training algorithm. Chapter~\ref{chap:experimental-protocol} defines the datasets, baselines, metrics, and implementation protocol. Chapter~\ref{chap:results-discussion} reports and interprets the deterministic and learned-ordering results. Finally, Chapter~\ref{chap:conclusion} summarizes the contributions, limitations, and future research directions.

%========================
% chapters/02_foundations.tex
%========================
\chapter{Foundations}
\label{chap:foundations}

This chapter introduces the concepts needed to understand the dissertation. The presentation is intentionally centered on the connection between graph topology and sequential generation, because the main object of study is not graph generation in the abstract, but the way in which a graph becomes a sequence before it is given to an autoregressive model.

\section{Graphs, Isomorphism, and Node Orderings}
\label{sec:graphs-orderings}

A graph is denoted by $G=(V,E)$, where $V$ is a finite set of nodes and $E\subseteq V\times V$ is a set of edges. In this dissertation, the experimental graphs are undirected and unweighted, although the sequential generation framework can be extended to richer settings. The number of nodes is $n=|V|$, and the number of edges is $m=|E|$. An adjacency matrix $A\in\{0,1\}^{n\times n}$ represents the same graph only after an ordering of the nodes has been chosen. If the nodes are permuted, the adjacency matrix changes even though the abstract graph remains the same.

This distinction is central. Two graphs $G=(V,E)$ and $G'=(V',E')$ are isomorphic when there exists a bijection $f:V\rightarrow V'$ such that $(u,v)\in E$ if and only if $(f(u),f(v))\in E'$. A generative model that learns graphs should ideally assign the same structural meaning to all isomorphic labelings. However, a neural model trained on matrices or sequences typically observes a particular labeling. Node ordering is therefore the bridge between the invariant object and the ordered representation used for learning.

For a graph $G$, let $\mathcal{O}_G$ be the set of all admissible node orderings. An ordering $\sigma=(v_1,\ldots,v_n)$ is a permutation of $V$. Once $\sigma$ is fixed, the graph can be translated into a sequence of construction actions. In DGMG, for example, the position of a node in $\sigma$ determines when it is introduced, which previous nodes can be selected as edge destinations, and when the model must decide to stop adding edges for the current node. Thus, ordering does not only relabel the graph; it changes the temporal structure of the supervised learning problem.

Ordering affects the generator through three mechanisms. The first is locality. If the correct edge destination is usually close in the ordering, then the generator can rely on recent construction context. The second is entropy. If an ordering makes one destination much more plausible than the others, the action distribution becomes sharper and easier to learn. The third is regularity. If similar graphs induce similar action patterns, training examples reinforce one another; if equivalent graphs induce highly variable sequences, the model must learn many incompatible histories.

These mechanisms explain why the dissertation does not treat graph ordering as a mere preprocessing operation. A fixed ordering rule already encodes a hypothesis about which partial graphs should be easy to continue. A learned ordering policy extends this idea by making the hypothesis trainable and graph-specific.

The dissertation studies several families of graphs because each one emphasizes a different structural constraint. Cycles require every node to have degree two and require a single ring-like closure. Trees require connectivity and acyclicity. Binary trees add a branching constraint. Stars are dominated by a single hub connected to all leaves. Barab\'asi--Albert graphs are sparse scale-free graphs generated by growth and preferential attachment \cite{barabasi1999emergence}. These families were chosen because their topology makes the effect of ordering interpretable.

\section{Autoregressive Graph Generation}
\label{sec:autoregressive-generation}

An autoregressive graph generator decomposes graph generation into a sequence of conditional decisions. Rather than sampling a whole graph at once, it samples actions $a_1,\ldots,a_L$ and defines
\begin{equation}
    p_{\phi}(a_1,\ldots,a_L)
    = \prod_{t=1}^{L} p_{\phi}(a_t \mid a_{<t}),
\end{equation}
where $\phi$ denotes the generator parameters and $a_{<t}$ is the history of previous actions. In graph generation, the history corresponds to a partial graph. This decomposition is attractive because it mirrors a constructive process: nodes and edges can be added incrementally, and validity constraints can be checked during generation.

The price of this decomposition is sensitivity to the chosen action sequence. If a graph can be serialized in multiple ways, each serialization produces a different factorization of the same underlying object. For a graph $G$ and ordering $\sigma$, let $\mathcal{T}(G,\sigma)$ denote the translator that produces a sequence $a=(a_1,\ldots,a_L)$. The negative log-likelihood optimized during teacher-forced training is
\begin{equation}
    \mathcal{L}_{\mathrm{NLL}}(\phi;G,\sigma)
    = -\sum_{t=1}^{L}\log p_{\phi}(a_t \mid a_{<t}),
    \qquad a=\mathcal{T}(G,\sigma).
\end{equation}
Changing $\sigma$ changes $a$ and therefore changes the conditional prediction task. In this sense, ordering is a structural inductive bias imposed before the generator is trained.

\section{Graph Neural Networks}
\label{sec:gnn-foundations}

Graph Neural Networks encode graph structure by repeatedly exchanging messages along edges \cite{scarselli2009gnn,gilmer2017mpnn}. A general message-passing layer updates a node representation by aggregating information from its neighborhood:
\begin{equation}
    h_v^{(k+1)} = \psi^{(k)}\left(h_v^{(k)},\; \square_{u\in\mathcal{N}(v)}\, \eta^{(k)}(h_v^{(k)},h_u^{(k)},e_{uv})\right),
\end{equation}
where $h_v^{(k)}$ is the representation of node $v$ at layer $k$, $\mathcal{N}(v)$ is the neighborhood of $v$, $\eta$ is a message function, $\square$ is a permutation-invariant aggregation operator such as sum or mean, and $\psi$ is an update function. Graph Convolutional Networks simplify this update through normalized neighborhood aggregation \cite{kipf2017gcn}, while Graph Attention Networks use learned attention weights to let the model assign different importance to different neighbors \cite{velickovic2018gat}.

GNNs are important in this dissertation for two reasons. First, autoregressive models uses graph representations to condition the next construction decision. Second, the learned-ordering policy must observe the entire input graph and choose a node to place next in the ordering. A policy that ignores the graph would reduce to a generic permutation sampler, whereas a GNN-based policy can condition its choices on degree, neighborhood structure, frontier effects, and global graph signals.

\section{Reinforcement Learning and Policy Gradients}
\label{sec:rl-foundations}

Reinforcement learning studies sequential decision problems in which an agent interacts with an environment and receives rewards \cite{sutton2018reinforcement}. In a Markov Decision Process, the agent observes a state $s_t$, samples an action $a_t$ from a policy $\pi_\theta(a_t\mid s_t)$, moves to a new state, and eventually receives a return. The objective is to maximize expected return:
\begin{equation}
    J(\theta)=\mathbb{E}_{\tau\sim\pi_\theta}[R(\tau)],
\end{equation}
where $\tau$ is a trajectory and $R(\tau)$ is its return.

When actions are discrete, the reward is delayed, and the environment is not differentiable, policy-gradient methods are a natural choice. The REINFORCE estimator \cite{williams1992reinforce} uses the log-derivative identity to write
\begin{equation}
    \nabla_\theta J(\theta)
    =\mathbb{E}_{\tau\sim\pi_\theta}\left[R(\tau)\nabla_\theta\log p_\theta(\tau)\right].
\end{equation}
In the node-ordering problem, the trajectory is the sequence of selected nodes, and the reward is computed only after the resulting ordering has been used to train and evaluate the graph generator. This makes the problem a direct instance of delayed-reward structured decision-making.

%\section{Why Ordering Changes the Learning Problem}
%\label{sec:ordering-learning-problem}

%Ordering affects the generator through three mechanisms. The first is locality. If the correct edge destination is usually close in the ordering, then the generator can rely on recent construction context. The second is entropy. If an ordering makes one destination much more plausible than the others, the action distribution becomes sharper and easier to learn. The third is regularity. If similar graphs induce similar action patterns, training examples reinforce one another; if equivalent graphs induce highly variable sequences, the model must learn many incompatible histories.

%These mechanisms explain why the dissertation does not treat graph ordering as a mere preprocessing operation. A fixed ordering rule already encodes a hypothesis about which partial graphs should be easy to continue. A learned ordering policy extends this idea by making the hypothesis trainable and graph-specific.

%========================
% chapters/03_related_work.tex
%========================
\chapter{Related Work}
\label{chap:related-work}

Graph generation has been studied from several perspectives, ranging from classical random graph models to modern neural generators. This chapter reviews the literature most relevant to the dissertation, with emphasis on methods that either rely on graph-to-sequence conversion or explicitly address the ordering problem.

\section{Graph Generative Models}
\label{sec:related-graph-generative-models}

Classical graph models define distributions through explicit structural assumptions. Erd\H{o}s--R\'enyi graphs model independent edge formation \cite{erdos1960evolution}, Watts--Strogatz graphs capture small-world behavior \cite{watts1998collective}, and Barab\'asi--Albert graphs model growth with preferential attachment \cite{barabasi1999emergence}. These models are useful because their assumptions are transparent, but they are limited when the target distribution contains complex motifs, constraints, attributes, or domain-specific generation rules.

Neural graph generation attempts to learn such distributions from data. Variational methods encode graphs into latent variables and decode them back into graphs \cite{kipf2016vgae,simonovsky2018graphvae}; adversarial methods train a generator against a discriminator \cite{goodfellow2014generative,bojchevski2018netgan}; flow-based models define invertible transformations for graph objects \cite{shi2020graphaf}; and diffusion models generate graphs through iterative denoising \cite{vignac2023digress,chen2023edge}. These approaches differ in how they handle permutation invariance, discrete structure, and validity constraints, but they share the general goal of learning a graph distribution without explicitly hand-designing a random graph model for each domain.

Autoregressive generators form a particularly relevant branch because they construct graphs through sequences of local decisions. DGMG generates graphs by repeatedly deciding whether to add a node, add an edge, and choose a destination \cite{li2018learning}. GraphRNN decomposes graph generation into a node sequence and edge sequences conditioned on previous nodes \cite{you2018graphrnn}. GRAN improves scalability with recurrent attention and block-wise generation \cite{liao2020gran}. Edge-based recurrent models also build graphs through sequential edge decisions \cite{bacciu2019edgegraphgen}. These models are expressive and naturally support constructive constraints, but they expose the ordering problem directly: before training, the graph must be represented as an ordered object.

\section{The Ordering Problem in Autoregressive Generation}
\label{sec:related-ordering-problem}

The dependence of autoregressive graph models on node order has motivated a growing body of work. Chen et al. formalize the role of node sequence in graph generation and show that order matters probabilistically \cite{chen2021ordermatters}. Bu et al. argue that ordering should not be treated as a secondary detail and analyze how different orderings affect autoregressive graph generation \cite{bu2023let}. Other approaches attempt to overcome ordering by marginalizing, augmenting, or modifying the factorization over orderings \cite{cohenkarlik2024olr,wang2025loarm}. Transformer-based graph generators also revisit sequence design, since attention can model long-range dependencies but still requires a representation of the graph as tokens or construction actions \cite{zhao2023graphgpt,chen2025g2pt}.

The present dissertation is aligned with this literature but takes a specific position. Instead of trying to remove the ordering dependency, it treats ordering as an optimization variable. This is motivated by the empirical observation that orderings can be beneficial when they match the topology of the target family. A BFS ordering can simplify tree generation, a DFS-like ordering can simplify cycles, and a degree ordering can simplify stars. In this view, the problem is not merely that orderings introduce nuisance variation; it is also that a good ordering can be a useful structural prior.

\section{Deterministic Ordering Heuristics}
\label{sec:related-heuristics}

Deterministic ordering heuristics are widely used because they are simple, reproducible, and inexpensive compared with model training. Breadth-first search and depth-first search are classical traversals \cite{cormen2009introduction}. Degeneracy ordering, also known as smallest-last ordering, repeatedly removes low-degree vertices and is connected to sparsity structure \cite{matula1983smallestlast}. Maximum cardinality search selects vertices according to how many already selected neighbors they have and is historically linked to sparse matrix factorization and chordal graph recognition \cite{rose1976algorithmic}. Spectral ordering uses eigenvectors of graph matrices to impose a global one-dimensional layout, drawing on spectral graph theory \cite{fiedler1973algebraic}.

These heuristics are not equivalent from the perspective of an autoregressive generator. A traversal ordering emphasizes reachability and local expansion; a degree ordering emphasizes hub exposure; degeneracy emphasizes sparse cores; and spectral ordering emphasizes global geometry. The dissertation uses these heuristics as controlled interventions: the generator architecture remains fixed, while the ordering changes. This makes it possible to isolate the effect of the graph-to-sequence map.

\section{Learning to Order}
\label{sec:related-learning-order}

Learning an ordering is more difficult than selecting a deterministic heuristic because the space of permutations is discrete and combinatorial. The order of $n$ nodes contains $n!$ possibilities, and the usefulness of a partial ordering may only become visible after the full graph has been translated, the generator has been trained, and samples have been evaluated. This makes the problem naturally related to sequential decision-making.

Reinforcement learning has been used for combinatorial optimization and sequencing tasks because it can optimize policies through delayed rewards \cite{sutton2018reinforcement}. Policy-gradient methods are especially appropriate when the action space is discrete and the reward cannot be differentiated with respect to the decisions \cite{williams1992reinforce}. In the context of this dissertation, selecting an ordering is treated as a graph-conditioned trajectory. The policy receives the graph and the set of already selected nodes, chooses the next node, and continues until all nodes have been selected. The reward is based on downstream DGMG behavior. This differs from supervised ordering prediction because there is no unique ground-truth ordering; the best ordering is defined by its effect on the generator.

The use of attention in the final policy architecture is also motivated by developments in neural sequence modeling and graph representation learning. Attention mechanisms allow a model to weight context elements differently \cite{vaswani2017attention}, and graph attention extends this idea to neighborhoods \cite{velickovic2018gat}. In the learned-ordering problem, attention is useful because different nodes may become relevant at different stages of the ordering trajectory. Early choices may benefit from hub information, while later choices may depend more strongly on frontier structure or remaining-node connectivity.

\section{Position of This Dissertation}
\label{sec:related-position}

The dissertation contributes to ordering-aware graph generation in three ways. First, it provides a topology-dependent empirical study of deterministic orderings under a fixed autoregressive model generator. Second, it formulates learned ordering as a graph-conditioned reinforcement-learning problem rather than as a supervised ranking problem. Third, it evaluates a stronger attention-based policy on Barab\'asi--Albert graphs and shows that learned ordering can improve the trade-off among size validity, connectivity, and hub-related statistics. The central perspective is that ordering should be studied as part of the model, because the sequence presented to an autoregressive generator is one of the main determinants of what the generator can learn efficiently.

%========================
% chapters/04_dgmg_reproduction.tex
%========================
\chapter{DGMG Reproduction and Graph-to-Sequence Translation}
\label{chap:dgmg-reproduction}

This chapter describes the reproduction of the Deep Generative Model of Graphs used as the experimental environment for the dissertation. The goal of the reproduction is not only to obtain a working graph generator, but also to expose the point at which node ordering enters the learning pipeline. Once this point is made explicit, the same generator can be trained under different graph-to-sequence maps and the effect of ordering can be measured directly.

\section{Constructive Generation in DGMG}
\label{sec:dgmg-constructive-generation}

DGMG is an autoregressive graph generator that constructs a graph through a sequence of decisions \cite{li2018learning}. At each stage, the model maintains a partial graph $G_t$ and predicts the next action conditioned on a learned representation of that partial graph. The action set includes adding a new node, deciding whether to add an edge from the current node, selecting the destination of that edge among previous nodes, and stopping edge addition for the current node. The process terminates when the model decides not to add another node.

The constructive nature of DGMG is important for this dissertation because it makes the ordering dependency transparent. During supervised training, a target graph must be converted into the exact sequence of actions that would construct it. This conversion requires a node ordering. If the ordering changes, the node-addition times change, the set of available edge destinations at each step changes, and the target sequence changes. The model architecture may be identical, but the training problem is different.

Formally, given a graph $G$ and an ordering $\sigma=(v_1,\ldots,v_n)$, the translator produces an action sequence
\begin{equation}
    a = \mathcal{T}(G,\sigma)=(a_1,\ldots,a_L).
\end{equation}
The generator is trained by maximizing the likelihood of this sequence or, equivalently, minimizing
\begin{equation}
    \mathcal{L}_{\mathrm{NLL}}(\phi;G,\sigma)
    = -\sum_{t=1}^{L}\log p_\phi(a_t\mid a_{<t}).
\end{equation}
This equation is the simplest expression of the ordering problem. The graph $G$ is fixed, but $\sigma$ changes the loss surface seen by the model.

\section{Graph Representation and Decision Modules}
\label{sec:dgmg-modules}

The reproduced implementation follows the modular structure of DGMG. A graph embedding module computes node and graph representations for the current partial graph. Message passing updates node states through local neighborhoods, and a readout function aggregates node information into a graph-level vector. Separate decision modules then predict the next action type. The AddNode module decides whether generation should introduce a new node or terminate. The AddEdge module decides whether the current node should receive another edge. The ChooseDest module selects the destination of the new edge among nodes already present in the partial graph.

The ChooseDest decision is the most directly affected by ordering. If the correct destination is a recent or structurally obvious previous node, the decision can be learned with lower ambiguity. If the correct destination is one among many unrelated previous nodes, the model must distribute probability across a larger and less coherent candidate set. This is why the dissertation repeatedly interprets ordering through locality and entropy: a good ordering makes the correct destinations easier to identify from the partial graph.

\section{Teacher-Forced Training}
\label{sec:dgmg-training}

The reproduction uses teacher-forced maximum likelihood training. For each target graph, the translator provides the correct sequence of AddNode, AddEdge, and ChooseDest actions. During training, the model conditions on the ground-truth history rather than on its own sampled mistakes. This stabilizes optimization and allows a clear comparison among orderings, because each ordering deterministically defines the target sequence used for training.

Let $\mathcal{D}$ be the training set and let $\sigma_G$ be the ordering assigned to graph $G$. The empirical training objective is
\begin{equation}
    \min_\phi \; \frac{1}{|\mathcal{D}|}\sum_{G\in\mathcal{D}}\mathcal{L}_{\mathrm{NLL}}(\phi;G,\sigma_G).
\end{equation}
In the deterministic experiments, $\sigma_G$ is produced by a fixed heuristic such as BFS or degree ordering. In the learned-ordering experiments, $\sigma_G$ is sampled from the policy $\pi_\theta(\cdot\mid G)$. This distinction is the bridge from reproduction to the reinforcement-learning method introduced later.

\section{Translation Rule}
\label{sec:dgmg-translation-rule}

The translation rule used in this dissertation follows the same constructive logic for all orderings. Nodes are introduced according to $\sigma$. When node $v_i$ is introduced, only the nodes $v_1,\ldots,v_{i-1}$ are available as destinations. For every edge between $v_i$ and a previous node, the translator emits an AddEdge action followed by a ChooseDest action. Once all backward edges of $v_i$ have been emitted, the translator emits the action that stops edge addition for the current node. After the last node has been processed, the translator emits the final stop action for node generation.

The length of the sequence depends on the number of nodes and edges. For a simple undirected graph, each node addition contributes node-level actions, and each edge contributes an edge decision and a destination decision. The exact length depends on the adopted implementation details, but the important point is that the edge-destination part of the sequence is strongly conditioned by ordering. The same edge may be predicted early or late, and its destination may compete with few or many previous nodes depending on where the incident vertices appear in the ordering.

\begin{algorithm}[H]
\caption{Graph-to-sequence translation under a node ordering}
\label{alg:translator}
\begin{algorithmic}[1]
\Require Graph $G=(V,E)$ and ordering $\sigma=(v_1,\ldots,v_n)$
\Ensure Action sequence $a$
\State Initialize $a\leftarrow [\,]$
\For{$i=1$ to $n$}
    \State Append the action that adds node $v_i$
    \For{each $j<i$ such that $(v_i,v_j)\in E$}
        \State Append an AddEdge action
        \State Append a ChooseDest action pointing to $v_j$
    \EndFor
    \State Append the action that stops edge addition for $v_i$
\EndFor
\State Append the final action that stops node generation
\State \Return $a$
\end{algorithmic}
\end{algorithm}

\section{Role of the Reproduction}
\label{sec:dgmg-role}

The reproduction serves three purposes in the dissertation. First, it provides a faithful sequential generator whose behavior can be compared under different orderings. Second, it defines the exact interface used by the reinforcement-learning policy: the policy only needs to output an ordering, and the existing translator converts that ordering into DGMG supervision. Third, it makes the analysis interpretable. When performance changes, the change can be attributed to the ordering-induced sequence rather than to a change in generator architecture.

This design choice is important methodologically. Many graph generation studies compare different architectures, losses, or datasets. Here, the architecture is held fixed as much as possible. The central experimental intervention is the order in which the graph is revealed to the autoregressive model. This isolates node ordering as an independent modeling variable and prepares the ground for learning that variable directly.

%========================
% chapters/05_deterministic_orderings.tex
%========================
\chapter{Deterministic Node Orderings}
\label{chap:deterministic-orderings}

Before learning orderings, it is necessary to understand what fixed orderings do to the sequential generation problem. This chapter describes the deterministic orderings used as baselines and explains why each one may help or harm DGMG depending on the target graph topology. The discussion analyzes each ordering through the lens of the action sequence it produces, rather than through purely structural graph properties. This perspective is deliberate: the generator never observes the original graph directly during training. Instead, it observes only the sequence of construction actions derived from that graph. Therefore, what matters about an ordering is not how it looks geometrically, but how easy or difficult it makes the resulting action sequence to learn.

\section{Traversal-Based Orderings}
\label{sec:traversal-orderings}

Breadth-first search (BFS) orders nodes by expanding a frontier level by level from a root. In sparse hierarchical graphs, this tends to place parents before children and keeps related nodes close in the sequence. For tree and binary-tree generation, this is a strong inductive bias: when a new node is introduced, its correct parent is usually located in the recent frontier, and the ChooseDest decision becomes comparatively simple. BFS can be weaker for cycles because expanding by layers does not necessarily preserve the ring as a nearly linear path, and it can expose multiple plausible backward destinations too early.

Depth-first search (DFS) follows a path as far as possible before backtracking. This makes it naturally suited to structures in which a long chain or near-chain is a good representation. Cycles are the clearest example in this dissertation. A DFS-like order can reveal almost the entire ring as a path and postpone closure to a small number of final decisions. The same behavior is less obviously optimal for trees, where excessive depth can separate siblings and make some branching patterns less regular, but DFS still preserves strong local continuity.

LexBFS is a lexicographic variant of breadth-first search. It uses labels induced by previously selected nodes to break ties in a way that preserves more structural information than a naive BFS traversal. In the experiments, LexBFS often behaves as a robust compromise: it retains frontier logic but can also create stable path-like orderings in some families. This explains why it performs strongly in cycles and trees.

\section{Centrality and Sparsity Orderings}
\label{sec:centrality-sparsity-orderings}

Degree ordering places high-degree nodes early. Its effect is particularly strong in hub-dominated graphs. In a star, selecting the hub first collapses most later destination choices onto a single node, making the sequence extremely regular. Degree ordering can also help in scale-free graphs, where early exposure of hubs may support the preferential-attachment structure. However, it can be poor for acyclic families such as trees if centrality does not correspond to a parent-child construction process.

Degeneracy ordering is based on the repeated removal of low-degree nodes. It is often associated with sparse graph structure and core decomposition. From the perspective of sequential generation, degeneracy can create orderings that reveal a sparse backbone while avoiding some of the instability of pure degree ordering. In Barab\'asi--Albert graphs, this makes degeneracy a strong compromise: it helps preserve size and connectivity while still respecting the presence of hubs.

Maximum cardinality search (MCS) selects nodes according to how many already selected neighbors they have. This makes the ordering sensitive to the growing selected set. It often favors vertices that are well connected to the current prefix and can therefore reduce ambiguity in edge decisions. In the experiments, MCS is especially relevant for Barab\'asi--Albert graphs because it balances connectivity with realistic degree structure.

\section{Spectral and Random Orderings}
\label{sec:spectral-random-orderings}

Spectral ordering uses eigenvectors of graph matrices to impose a global one-dimensional layout. It is attractive because it can reflect global connectivity patterns, but the experiments show that global geometric structure is not always aligned with the local conditional decisions required by DGMG. A spectral layout may be graph-theoretically meaningful and still produce action sequences in which edge destinations are difficult to predict.

Random ordering is included as a structure-agnostic baseline. It intentionally removes the structural prior encoded by the other heuristics. If random ordering performed close to the best deterministic orderings, the dissertation's main premise would be weakened. Instead, the experiments show large gaps in several families, confirming that the graph-to-sequence map is an important modeling decision.

\section{Ordering as an Entropy-Shaping Mechanism}
\label{sec:ordering-entropy}

The main effect of deterministic ordering can be summarized as entropy shaping. In DGMG, the ChooseDest module must select one destination among previous nodes. When the correct destination is structurally obvious, the conditional distribution is easier to learn; when many previous nodes are plausible, the distribution becomes more diffuse. A good ordering reduces the effective ambiguity of these decisions by controlling which nodes are available and how they relate to the current node.

This explains the topology-dependent behavior observed later. Cycles favor orderings that make the graph look like a path until the closure decision. Trees favor orderings that expose parents before children in compact frontiers. Stars favor orderings that reveal the hub immediately. Barab\'asi--Albert graphs do not have a single dominant requirement; they combine hub exposure, sparse connectivity, and size control. As a result, different deterministic orderings optimize different metrics.

The deterministic study motivates learned ordering for two reasons. First, no fixed ordering is universally best. A heuristic that is excellent for one topology can fail for another. Second, even within a single family, the best behavior may be hybrid. For example, a Barab\'asi--Albert graph may benefit from degree-like choices early in the sequence to expose hubs, but from BFS-like or MCS-like choices later to preserve connectivity. A handcrafted rule must commit to one principle, whereas a learned policy can condition its choices on the graph and on the partial ordering already built.

The next chapter turns this motivation into a formal reinforcement-learning problem. The policy will not be trained to imitate a known ordering. Instead, it will be trained to produce orderings that improve the downstream generator.

%========================
% chapters/06_rl_ordering.tex
%========================
\chapter{Learning Node Orderings with Deep Reinforcement Learning}
\label{chap:rl-ordering}

This chapter presents the main methodological contribution of the dissertation. The deterministic study shows that node ordering can strongly change DGMG performance and that different topologies favor different orderings. The natural next step is to remove the requirement that the ordering be fixed by hand. The proposed method treats node ordering as a graph-conditioned sequential decision problem and trains a policy to construct orderings that improve the downstream graph generator.

\section{Problem Formulation}
\label{sec:rl-problem-formulation}

Let $\mathcal{D}$ be a dataset of graphs. For each graph $G\in\mathcal{D}$, the set of possible orderings is $\mathcal{O}_G$. A deterministic heuristic selects a single ordering $\sigma=h(G)$. In contrast, the learned method defines a stochastic policy $\pi_\theta(\sigma\mid G)$ over complete orderings. Once an ordering is sampled, the translator produces an action sequence $\mathcal{T}(G,\sigma)$, and DGMG is trained on that sequence.

The ideal objective would select policy parameters $\theta$ that minimize downstream generation loss after the generator has been trained under the policy-induced sequences. This yields a bilevel view of the problem. Let $\phi$ denote the generator parameters. In the lower-level task, for a given policy $\pi_\theta$, we seek the optimal generator parameters $\phi^\star(\theta)$ that minimize the expected training loss over the distribution of orderings:
\begin{equation}
    \phi^\star(\theta) = \arg\min_{\phi} \;\mathbb{E}_{G\sim\mathcal{D}}\left[\mathbb{E}_{\sigma\sim\pi_\theta(\cdot\mid G)}\left[\mathcal{L}_{\mathrm{NLL}}\!\left(\phi,\mathcal{T}(G,\sigma)\right)\right]\right].
    \label{eq:bilevel-phi}
\end{equation}
The optimal policy parameters $\theta^\star$ are found by solving the upper-level task to maximize the expected downstream reward $R$, evaluated at the lower-level equilibrium $\phi^\star(\theta)$:
\begin{equation}
    \theta^\star = \arg\max_{\theta} \;\mathbb{E}_{G\sim\mathcal{D}}\left[\mathbb{E}_{\sigma\sim\pi_\theta(\cdot\mid G)}\left[R\bigl(G,\sigma,\phi^\star(\theta)\bigr)\right]\right],
    \label{eq:bilevel-theta}
\end{equation}
where $R(G,\sigma,\phi)$ is a delayed reward that evaluates the ordering $\sigma$ against the current generator state $\phi$.

Solving Eq.~\eqref{eq:bilevel-phi} to convergence for every update of $\theta$ is computationally prohibitive. Therefore, we adopt an alternating optimization strategy, detailed in Section~\ref{sec:rl-alternating-optimization}, in which the generator and the ordering policy are updated in interleaved stages, effectively performing stochastic gradient steps on both levels.

This design reflects the fact that the ordering is valuable only through its effect on the generator. There is no ground-truth label saying which ordering is correct. An ordering is good if it makes the sequential learning problem easier and leads to better samples.

\section{Markov Decision Process}
\label{sec:rl-mdp}

The ordering process is formulated as a finite-horizon Markov Decision Process defined
through the tuple $(\mathcal{S}, \mathcal{A}, \mathcal{P}, \mathcal{R}, \gamma)$, where
$\mathcal{S}$ is the state space, $\mathcal{A}$ is the action space, $\mathcal{P}$ is the
state transition function, $\mathcal{R}$ is the reward function, and $\gamma \in [0,1]$
is the discount factor that controls how much future rewards are weighted relative to
immediate ones. Since the task is episodic and finite, we assume $\gamma = 1$, meaning
all steps in the trajectory contribute equally to the cumulative reward, keeping the
objective focused on the quality of the complete ordering rather than on any intermediate
decision in isolation.

At the beginning of an episode, the input graph $G=(V,E)$ is fixed and the selected set is empty. At step $t$, the state is
\begin{equation}
    s_t=(G,S_t),
\end{equation}
where $S_t\subseteq V$ is the set of nodes already selected. The action space is the set of remaining nodes,
\begin{equation}
    \mathcal{A}(s_t)=V\setminus S_t.
\end{equation}
After the policy selects $u_t\in\mathcal{A}(s_t)$, the state transitions deterministically to $S_{t+1}=S_t\cup\{u_t\}$. The episode ends after $n$ steps, producing the ordering $\sigma=(u_1,\ldots,u_n)$.

The probability of a complete ordering under the policy is
\begin{equation}
    \pi_\theta(\sigma\mid G)
    = \prod_{t=1}^{n}\pi_\theta(u_t\mid G,S_t),
\end{equation}
where invalid actions are masked so that already selected nodes cannot be chosen again.

The reward is defined at the trajectory level. An episode terminates once all nodes have been selected. At termination, we obtain a valid node permutation $\sigma=(a_1,\ldots,a_n)$ defined by the sequence of actions in $\tau$. Let $\phi$ denote the current DGMG parameters after a generator-update stage using sequences induced by the policy. The step-wise reward component is defined as
\begin{equation}
    \mathcal{R}(s_t,a_t) = \log p_\phi(a_t \mid a_{<t}),
\end{equation}
noting that $s_t$ implicitly contains the information of the previous decisions $a_{<t}$.

A central design principle of the method is that an ordering should not be rewarded in isolation. Instead, the reward must reflect the quality of the downstream generator trained under that ordering. The trajectory return $R(\tau)$ is therefore defined as a combination of predictive fit $R_{\mathrm{pred}}(\tau)$ and structural validity $R_{\mathrm{struct}}(\tau)$. This avoids optimizing the policy for spurious sequence patterns by tying the objective to both likelihood and graph-level structural fidelity.

The predictive term is defined as follows. Given a graph $G$ and a permutation $\sigma$, the validation loss is
\begin{equation}
    \mathcal{L}_{\mathrm{val}}(\phi,\mathcal{T}(G,\sigma))
    = -\sum_{t=1}^{n}\log p_\phi(a_t\mid a_{<t}).
\end{equation}
Let $\mathcal{D}_{\mathrm{val}}=\{G^{(i)}\}_{i=1}^{|\mathcal{D}_{\mathrm{val}}|}$ be the validation set. The average validation loss is
\begin{equation}
    \overline{\mathcal{L}}_{\mathrm{val}}(\tau)
    = \frac{1}{|\mathcal{D}_{\mathrm{val}}|}\sum_{i=1}^{|\mathcal{D}_{\mathrm{val}}|}\mathcal{L}_{\mathrm{val}}\!\left(\phi,\mathcal{T}(G^{(i)},\sigma)\right).
\end{equation}
Since lower values of $\overline{\mathcal{L}}_{\mathrm{val}}$ are better, the predictive reward is its negative value:
\begin{equation}
    R_{\mathrm{pred}}(\tau)
    = -\overline{\mathcal{L}}_{\mathrm{val}}(\tau)
    = \frac{1}{|\mathcal{D}_{\mathrm{val}}|}\sum_{i=1}^{|\mathcal{D}_{\mathrm{val}}|}\sum_{t=1}^{n}\log p_\phi(a_t\mid a_{<t}).
\end{equation}

The structural term $R_{\mathrm{struct}}(\tau)$ is tailored to the graph family of interest. To estimate it, candidate graphs are sampled from the current DGMG model and evaluated using a family-specific validity predicate. In the Barab\'asi--Albert setting, this structural component incorporates size control, connectivity, and distributional statistics related to scale-free behavior, including degree inequality and hub concentration. In this way, the policy is encouraged not only to improve predictive fit, but also to induce orderings that support the recovery of the characteristic heavy-tailed structure of BA graphs.

The final composite reward is
\begin{equation}
    R(\tau) = R_{\mathrm{pred}}(\tau) + \lambda\,R_{\mathrm{struct}}(\tau),
    \label{eq:reward-composite}
\end{equation}
where $\lambda>0$ controls the trade-off between predictive fit and structural validity. To make the dependency explicit, we write $R(\tau)\equiv R(G,\sigma,\phi)$, since the trajectory $\tau$ defines a node permutation $\sigma$ whose quality depends on the input graph $G$ and on the generator parameters $\phi$ used to evaluate it.

The reward is delayed and assigned only after the complete ordering has been translated and evaluated through the generator. This delayed reward is the main reason for using policy gradients rather than ordinary backpropagation through the ordering decisions.

\section{Policy Architecture}
\label{sec:rl-policy-architecture}

The policy must choose a node while considering both the original graph and the partial ordering. The architecture therefore computes node embeddings with graph message passing and augments them with a selected-node indicator. In the final method, the encoder uses multiple graph layers and an attention mechanism. The additional layers increase the receptive field of each node, while attention allows the policy to weight neighboring signals differently depending on the current graph structure. This makes the architecture especially suited to graphs with heterogeneous local influence patterns, such as Barab\'asi--Albert networks with hubs and long-tailed degree distributions.

Let $x_v$ be the input feature vector for node $v$. In the synthetic setting, node features are lightweight and may include structural descriptors such as normalized degree together with a binary indicator showing whether the node has already been selected. Given state $s_t=(G,S_t)$, the graph encoder computes node embeddings
\begin{equation}
    h_v^{(T_\sigma)} = \mathrm{GNN}_\theta(G,S_t,x_v),
\end{equation}
where $T_\sigma$ denotes the number of message-passing rounds used by the policy network. A graph-level context vector $h_G$ is computed by pooling node embeddings via a permutation-invariant readout:
\begin{equation}
    h_G^{(T_\sigma)} = \mathrm{READOUT}\bigl(\{h_v^{(T_\sigma)}\}_{v\in V}\bigr).
\end{equation}
The policy score for an available node $v$ is produced by an MLP:
\begin{equation}
    z_v^{(t)} = f_\theta\!\left([h_v^{(T_\sigma)}, h_G^{(T_\sigma)}, \mathbb{I}\{v\in S_t\}]\right).
\end{equation}
The distribution over next nodes is obtained by applying a masked softmax:
\begin{equation}
    \pi_\theta(v\mid G,S_t)
    = \frac{\exp(z_v^{(t)})}{\sum_{u\in V\setminus S_t}\exp(z_u^{(t)})},
    \qquad v\in V\setminus S_t.
\end{equation}
The mask enforces validity of the permutation: without it, the policy could repeatedly select the same node, which would not define a valid ordering. The policy is sampled stochastically during training to preserve exploration and executed greedily during evaluation.

The final version of the method introduces a degree-based warm start. In the first iteration, the ordering follows degree information, giving the generator and policy a stable initial structural prior. This is particularly useful for Barab\'asi--Albert graphs, where hubs are central to the data-generating mechanism. After this warm start, the policy is free to adapt and can learn deviations from pure degree ordering when they improve validation loss or structural metrics.

\section{Policy-Gradient Objective}
\label{sec:rl-policy-gradient}

Because the ordering is discrete, gradients from the generator loss cannot be directly propagated through the selected node indices. The policy is therefore optimized with a likelihood-ratio estimator. For a fixed graph $G$, the expected reward is
\begin{equation}
    J_G(\theta)=\mathbb{E}_{\sigma\sim\pi_\theta(\cdot\mid G)}[R(G,\sigma)].
\end{equation}
Averaging over the training distribution gives
\begin{equation}
    J(\theta)=\mathbb{E}_{G\sim\mathcal{D}}\left[J_G(\theta)\right].
\end{equation}
The following theorem formalizes the gradient used by the method.

\begin{theorem}[Policy gradient for learned node ordering]
\label{thm:policy-gradient-ordering}
Let $G=(V,E)$ be a finite graph and let $\mathcal{O}_G$ be the set of all valid node orderings. Suppose that $\pi_\theta(\sigma\mid G)>0$ is differentiable with respect to $\theta$ for every $\sigma\in\mathcal{O}_G$, and let $R(G,\sigma)$ be a bounded reward assigned to a complete ordering. Then
\begin{equation}
    \nabla_\theta J_G(\theta)
    = \mathbb{E}_{\sigma\sim\pi_\theta(\cdot\mid G)}
    \left[ R(G,\sigma)\nabla_\theta\log \pi_\theta(\sigma\mid G) \right].
\end{equation}
If the ordering is sampled sequentially as $\sigma=(u_1,\ldots,u_n)$, then
\begin{equation}
    \nabla_\theta\log \pi_\theta(\sigma\mid G)
    = \sum_{t=1}^{n}\nabla_\theta\log \pi_\theta(u_t\mid G,S_t).
\end{equation}
\end{theorem}

\begin{proof}
By definition,
\begin{equation}
    J_G(\theta)=\sum_{\sigma\in\mathcal{O}_G}\pi_\theta(\sigma\mid G)R(G,\sigma).
\end{equation}
Since $\mathcal{O}_G$ is finite and $R$ is bounded, the gradient can be moved inside the sum:
\begin{equation}
    \nabla_\theta J_G(\theta)
    =\sum_{\sigma\in\mathcal{O}_G}\nabla_\theta\pi_\theta(\sigma\mid G)R(G,\sigma).
\end{equation}
Using the identity $\nabla_\theta\pi_\theta(\sigma\mid G)=\pi_\theta(\sigma\mid G)\nabla_\theta\log\pi_\theta(\sigma\mid G)$ gives
\begin{equation}
    \nabla_\theta J_G(\theta)
    =\sum_{\sigma\in\mathcal{O}_G}\pi_\theta(\sigma\mid G)R(G,\sigma)\nabla_\theta\log\pi_\theta(\sigma\mid G),
\end{equation}
which is the stated expectation. The sequential factorization follows from
\begin{equation}
    \pi_\theta(\sigma\mid G)=\prod_{t=1}^{n}\pi_\theta(u_t\mid G,S_t),
\end{equation}
so the logarithm of the trajectory probability is the sum of the logarithms of the step probabilities. Differentiating this sum completes the proof.
\end{proof}

In practice, a baseline can be subtracted from the reward to reduce variance without changing the expected gradient:
\begin{equation}
    \nabla_\theta J_G(\theta)
    = \mathbb{E}\left[(R(G,\sigma)-b(G))\nabla_\theta\log\pi_\theta(\sigma\mid G)\right].
\end{equation}
This is useful because the reward is noisy: it depends on generator training, validation loss, and sample-based structural metrics. An entropy regularizer may also be added to encourage exploration and prevent premature policy collapse:
\begin{equation}
    J_{\mathrm{ent}}(\theta)
    = J(\theta) + \beta\,\mathbb{E}\!\left[\sum_{t=1}^{n}\mathcal{H}\bigl(\pi_\theta(\cdot\mid s_t)\bigr)\right],
\end{equation}
where $\beta\ge0$ controls the strength of the entropy bonus.

\section{Alternating Optimization}
\label{sec:rl-alternating-optimization}

The learned-ordering method alternates between updating the generator and updating the policy. At the beginning of an outer iteration, the policy assigns an ordering to each training graph. The translator converts each ordered graph into a DGMG action sequence. The generator is then trained for a fixed number of epochs on these sequences. After this generator update, validation loss and sample metrics are computed. These values define the reward used to update the policy.

The outer loop is essential. If the policy were updated after a single sequence without retraining the generator, the reward would mostly reflect the current generator's limitations rather than the usefulness of the ordering. If the generator were trained to convergence after each policy change, the method would become computationally impractical. The alternating procedure is a compromise that lets the generator and policy co-adapt.

\begin{algorithm}[H]
\caption{Attention RL ordering with alternating DGMG training}
\label{alg:attention-rl-ordering}
\begin{algorithmic}[1]
\Require Training graphs $\mathcal{D}_{\mathrm{train}}$, validation graphs $\mathcal{D}_{\mathrm{val}}$, policy $\pi_\theta$, generator $p_\phi$
\For{outer iteration $i=1,\ldots,I$}
    \For{\textbf{each} graph $G\in\mathcal{D}_{\mathrm{train}}$}
        \If{$i=1$}
            \State Compute degree ordering $\sigma$
        \Else
            \State Sample or compute ordering $\sigma\sim\pi_\theta(\cdot\mid G)$
        \EndIf
        \State Translate graph into action sequence $a=\mathcal{T}(G,\sigma)$
    \EndFor
    \State Update DGMG parameters $\phi$ by minimizing teacher-forced NLL on the translated sequences
    \State Evaluate validation NLL and sample graphs from $p_\phi$
    \State Compute structural metrics and the reward $R(\tau)$ from Eq.~\eqref{eq:reward-composite}
    \For{$k=1,\ldots,K$ policy steps}
        \State Estimate policy gradient with REINFORCE
        \State Update $\theta$
    \EndFor
\EndFor
\State \Return trained policy $\pi_\theta$ and generator $p_\phi$
\end{algorithmic}
\end{algorithm}

The variable $i$ in the outer loop is not the same as the position $t$ in the ordering trajectory. The trajectory index $t$ runs from $1$ to $n$ within a single graph and selects the nodes of one ordering. The outer-loop index $i$ counts rounds of joint optimization between the ordering policy and the generator. This distinction is important when interpreting the algorithm: a single outer iteration contains many ordering decisions, one sequence per graph, and one generator-training stage. The parameter $I$ controls how many times generator learning and policy learning alternate, while $K$ controls how strongly the policy is updated within each alternation.

\section{Why Attention RL Ordering Differs from a Fixed Heuristic}
\label{sec:rl-hybrid-behavior}

A deterministic heuristic applies the same rule at every step. Degree ordering always prioritizes degree, BFS always follows frontier expansion, and MCS always prioritizes connectivity to the selected set. A learned policy can imitate any of these behaviors when useful, but it can also switch behavior during the ordering. For example, in a Barab\'asi--Albert graph, the policy may begin by selecting high-degree nodes to reveal hubs and then select nodes that preserve connectivity or complete local neighborhoods. This hybrid possibility is the main reason for learning an ordering rather than merely selecting the best deterministic baseline.

The attention mechanism strengthens this ability because the policy does not rely on a single scalar score such as degree. It can combine local neighborhoods, global graph context, and the current selected set. As the ordering progresses, the selected-node mask changes the state, and the policy distribution can change accordingly. The ordering is therefore not just graph-conditioned; it is history-conditioned. The three-layer message-passing stack further enlarges the effective receptive field before each selection step, and the non-random degree-based warm start ensures that the controller has a strong hub-aware prior from the very beginning of training, rather than having to rediscover hub structure through exploration.

%\section{Chapter Summary}
%\label{sec:rl-summary}

%This chapter formulated node ordering as a graph-conditioned reinforcement-learning problem. A complete ordering is treated as a trajectory, each action selects one unselected node, and the reward is determined by downstream DGMG performance. The bilevel structure makes explicit why this is harder than standard supervised learning: the reward is only defined after the generator has been updated, and there is no ground-truth ordering to imitate. The policy-gradient theorem justifies the REINFORCE update used for the discrete ordering decisions, and the alternating algorithm couples policy improvement with generator training. The final Attention RL ordering method adds graph attention, multiple message-passing layers, and a degree warm start, producing the version evaluated in the final Barab\'asi--Albert experiments.

%========================
% chapters/07_experimental_protocol.tex
%========================
\chapter{Experimental Protocol}
\label{chap:experimental-protocol}

This chapter describes the datasets, baselines, metrics, and implementation protocol used to evaluate deterministic and learned node orderings. The experiments are designed to isolate the effect of ordering while keeping the downstream generator fixed. The central question is whether the sequence induced by an ordering makes an autoregressive model algorithm better or worse at learning the target graph family.

\section{Graph Families}
\label{sec:protocol-graph-families}

The deterministic study uses five synthetic graph families. Cycles are included because they require a clean ring structure and make the effect of path-like orderings easy to observe. Trees are included because they test whether an ordering preserves parent-child locality and acyclicity. Binary trees add a stricter branching constraint, making them a more demanding version of the tree setting. Stars isolate the effect of hub exposure: once the hub appears early, most destination decisions become simple. Barab\'asi--Albert graphs are included because they are less rigid than the other families and combine growth, sparse connectivity, and heavy-tailed degree structure.

The final learned-ordering experiments focus on Barab\'asi--Albert (BA) graphs. This choice is motivated by the fact that BA graphs expose a real trade-off among metrics. A model can generate graphs with good size but weak connectivity, or graphs with strong connectivity but weaker hub statistics. This makes the family more informative for learned ordering than a setting in which one exact validity predicate dominates the evaluation. The dataset is generated with the standard BA preferential-attachment mechanism using attachment parameter $m_{\mathrm{BA}}=2$ and target size range $N\in[20,60]$. It contains 6000 training graphs, 750 validation graphs, and 750 test graphs.

\section{Ordering Baselines}
\label{sec:protocol-baselines}

The deterministic baselines are degree ordering, BFS, DFS, degeneracy, LexBFS, MCS, spectral ordering, and random ordering. Each baseline is applied before translation, and the resulting action sequences are used to train the same DGMG architecture. The learned method is compared against these baselines under the same BA protocol. This comparison is intentionally strict: the only intended change is the mechanism that chooses the ordering.

Two learned-ordering configurations are evaluated. The first is an RL policy parameterized by a Graph Convolutional Network (GCN) \cite{kipf2017gcn}, which serves as a direct predecessor and establishes that the general learned-ordering framework is viable. The second and final configuration is referred to as Attention RL ordering. It differs from the earlier GCN policy by using an attention-based graph layer, increasing the policy depth to three graph layers, and initializing the first iteration with degree ordering. %These changes were introduced after observing that pure exploration in the space of permutations is unstable, especially in graphs whose structural identity depends strongly on hubs.

\section{Training and Evaluation}
\label{sec:protocol-training-evaluation}

For each deterministic baseline, the training graphs are first ordered by the corresponding heuristic. The graph-to-sequence translator then creates action sequences, and DGMG is trained with teacher-forced negative log-likelihood. Validation NLL is monitored during training, and generated samples are evaluated using structural metrics. This protocol ensures that differences in the results can be attributed to the ordering-induced sequences rather than to changes in model capacity.

For the GCN RL policy, DGMG is coupled with the learned ordering described in Chapter~\ref{chap:rl-ordering}, with the generator kept fixed at hidden dimension $H=16$, message-passing rounds $T=2$, learning rate $\alpha=10^{-4}$, batch size 1, and 25 training epochs per update stage.

For Attention RL ordering, training uses the alternating procedure described in Chapter~\ref{chap:rl-ordering}. The GCN-style encoder is replaced by an attention-based graph encoder implemented in PyTorch, using transformer-style attention to let the controller weigh structural cues more adaptively when selecting the next node in the ordering. The policy depth is increased to three graph layers, and the first outer iteration is warm-started from the degree ordering. Each subsequent outer iteration generates orderings, updates DGMG, evaluates the generator, and updates the policy. Validation and test losses are tracked across outer iterations, and the best checkpoint is selected according to a convergence criterion based on the difference between consecutive loss values across outer iterations.

\section{Metrics}
\label{sec:protocol-metrics}

The metrics depend on the graph family. For cycles, the main metric is the cycle ratio, which measures the fraction of generated graphs satisfying the cycle predicate. For trees, the tree ratio measures acyclicity and connectivity. For binary trees, the binary-tree ratio additionally checks the branching constraint. For stars, the star ratio measures whether the generated graph has the expected single-hub structure. In all these families, valid-size ratio is also reported because a graph with the correct structure but incorrect size range is not considered fully satisfactory.

For Barab\'asi--Albert graphs, evaluation is necessarily multi-objective. The valid-size ratio measures whether generated graphs fall in the expected node-count range. The connected ratio measures whether generated graphs are connected. The average size and average number of edges capture whether the model reproduces the scale of the target distribution. The degree Gini coefficient measures inequality in the degree distribution, with higher values indicating stronger concentration of degree mass in a subset of nodes. Hubiness measures how dominant the largest hub is relative to the rest of the graph. These metrics are interpreted together because BA graphs do not have a single exact predicate that fully defines validity. Validation NLL is also monitored during training to ensure that structural metrics are supported by evidence of predictive fit.

\section{Computational Setup}
\label{sec:protocol-computational-setup}

All experiments were conducted on an Apple MacBook Pro equipped with an Apple M1 chip (2020), 8~GB of unified memory, and running macOS. The M1 chip integrates an 8-core CPU (four performance cores and four efficiency cores) and a 7-core GPU, with a 16-core Neural Engine. Although not a dedicated high-performance computing platform, the M1 architecture's unified memory design provides efficient data movement between the CPU and the on-chip accelerator, which is particularly relevant for the iterative alternating-optimization procedure that repeatedly loads model states and evaluates structural metrics.

The deterministic ordering experiments are computationally inexpensive, as each baseline requires only a single preprocessing pass followed by standard teacher-forced training. The learned-ordering procedure is substantially more demanding because it interleaves repeated policy updates with full DGMG retraining stages. For the Attention RL ordering run on Barab\'asi--Albert graphs, where the best validation checkpoint was reached at outer iteration 61, the complete alternating optimization procedure required approximately 7 hours of wall-clock time on the described hardware. This cost reflects the depth of the alternating loop, the size of the BA training set, and the additional overhead of sampling generated graphs for structural evaluation at each outer iteration.

These timing figures should be interpreted as indicative of the computational profile of the method rather than as absolute benchmarks, since runtimes depend on software implementation, dataset size, and hardware configuration. The purpose of reporting them is to situate the practical cost of treating node ordering as an optimization target relative to the negligible cost of applying a deterministic heuristic.

\begin{figure}[H]
    \centering
    \maybeincludegraphics[width=0.72\linewidth]{figures/rl_entropy_mean.png}
    \caption{Policy-entropy diagnostic across RL outer iterations, showing the mean entropy of the policy distribution over unselected nodes. The decreasing trend indicates that the policy becomes increasingly concentrated over the course of training, which is consistent with the transition from wide exploration to more confident node selection.}
    \label{fig:rl-entropy-mean}
\end{figure}

\section{Reproducibility Considerations}
\label{sec:protocol-reproducibility}

The experiments rely on stochastic training, stochastic sampling, and in the learned-ordering case stochastic policy trajectories. For this reason, the results should be interpreted through repeated evaluation on fresh samples rather than through a single generated graph. The BA results are reported on freshly sampled generated graphs, which is important because validation loss alone does not guarantee structural fidelity. The protocol therefore combines likelihood-based evidence with sample-based structural evidence.

A second reproducibility concern is the cost of the alternating RL procedure. Deterministic orderings are inexpensive preprocessing steps, while learned ordering requires repeated policy and generator updates. The purpose of the dissertation is not to claim that learned ordering is always cheaper than heuristics. Rather, it shows that ordering is sufficiently important to justify treating it as an object of optimization, especially when handcrafted heuristics exhibit trade-offs or fail to dominate across metrics. The source code for the experiments is publicly available at \url{https://github.com/prunebr/learning_to_order}.

%========================
% chapters/08_results_discussion.tex
%========================
\chapter{Results and Discussion}
\label{chap:results-discussion}

This chapter presents the empirical results and connects them to the central thesis of the dissertation: node ordering changes the effective learning problem faced by an autoregressive graph generator. The discussion first summarizes the deterministic-ordering study, then focuses on the final Attention RL ordering results on Barab\'asi--Albert graphs.

\section{Deterministic Ordering Results}
\label{sec:results-deterministic}

The deterministic experiments show that there is no universally optimal ordering. The best ordering depends on the target topology and on the structural property being measured. This result is important because it justifies the transition from fixed heuristics to learned ordering. If a single heuristic dominated all graph families, a learned policy would be less compelling. Instead, the empirical evidence shows that each ordering acts as a different structural bias.

Cycles are best served by orderings that make the graph appear as a nearly linear chain before closure. LexBFS and DFS achieve the strongest cycle reconstruction, with LexBFS reaching a cycle ratio of $0.9962$ and a valid ratio of $0.7179$. The interpretation is direct: when the ring is exposed as a path, most edge decisions are local, and the final closure is isolated to a small number of decisions. BFS, despite being powerful for trees, performs poorly in cycles because level-wise expansion does not preserve the sequential geometry of the ring.

Trees and binary trees favor BFS. In generic trees, BFS reaches a tree ratio of $0.9818$ and a valid ratio of $0.9790$. In binary trees, BFS reaches a binary-tree ratio of $0.9940$ and a valid ratio of $0.9915$. This confirms the parent-frontier interpretation: level-wise ordering tends to place the correct parent in a compact and meaningful candidate set. DFS, LexBFS, MCS, and degeneracy also perform well in these families, but BFS is the most aligned with the constructive process required by the generator.

Stars invert this pattern. The best ordering is degree ordering, which achieves a star ratio of $0.9980$ and a valid ratio of $0.9980$. In a star, the central decision is whether the hub appears early. Once the hub is introduced, almost every later edge decision points to it, and the sequence becomes extremely regular. This result shows that centrality-based orderings are not generally good or bad; they are good when the topology is hub-dominated.

Barab\'asi--Albert graphs are more nuanced. Table~\ref{tab:ba-deterministic-results} shows that different deterministic orderings occupy different points in the metric space. BFS obtains perfect connectivity but produces smaller graphs on average, with an average size of only $27.3$ nodes compared to the training range of $[20,60]$. Degree ordering gives perfect valid-size ratio and the largest average size, but connectivity drops sharply to $0.6490$. Degeneracy and MCS provide strong compromises among valid size, connectivity, and hub-related statistics. This multi-objective behavior is exactly the setting in which a learned policy may be useful.

\begin{table}[H]
\centering
\footnotesize
\begin{tabular}{lrrrrr}
\toprule
\textbf{Ordering} & \textbf{Avg. size} & \textbf{Valid size} & \textbf{Connected} & \textbf{Degree Gini} & \textbf{Hubiness} \\
\midrule
BFS        & 27.3110 & 0.8130 & \textbf{1.0000} & 0.2923 & 0.4282 \\
Degeneracy & 37.0770 & 0.9670 & 0.9740 & 0.3287 & 0.3993 \\
MCS        & 40.5150 & 0.9670 & 0.9110 & 0.3323 & 0.3886 \\
DFS        & 40.8660 & 0.9630 & 0.8280 & 0.3411 & 0.4202 \\
Random     & 46.2100 & 0.9940 & 0.7930 & 0.3328 & 0.3460 \\
Degree     & \textbf{50.3650} & \textbf{1.0000} & 0.6490 & \textbf{0.3516} & 0.3326 \\
Spectral   & 48.4210 & 0.9910 & 0.6010 & 0.3184 & 0.3085 \\
LexBFS     & 50.2580 & 0.9540 & 0.4910 & 0.3358 & 0.3237 \\
\bottomrule
\end{tabular}
\caption{Deterministic-ordering results on Barab\'asi--Albert graphs.}
\label{tab:ba-deterministic-results}
\end{table}

The deterministic results support an entropy-based interpretation of ordering. A good ordering reduces the uncertainty of the next action by placing structurally relevant nodes in the prefix. The relevant notion of ``structurally relevant'' changes with the graph family. In cycles, it means path continuity; in trees, it means parent-frontier locality; in stars, it means early hub exposure; in BA graphs, it means a balance among hubs, connectivity, and sparsity.

\section{Learned Ordering: GCN RL Baseline}
\label{sec:results-gcn-rl}

Before evaluating the full Attention RL ordering, we first examined a GCN-based RL policy as an intermediate step. In this configuration, the ordering policy is parameterized by a standard Graph Convolutional Network \cite{kipf2017gcn} and trained using the same alternating REINFORCE procedure described in Chapter~\ref{chap:rl-ordering}. The GCN RL ordering achieves an average generated size of $39.3220$ and a perfect valid-size ratio of $1.0000$, demonstrating that the general learned-ordering framework is viable and that size control is achievable. However, connectivity remains at $0.6334$, which is clearly below the strongest deterministic baselines and does not substantially improve over the degree ordering warm start. This result establishes the limitation that motivates the attention-based upgrade: while a learned policy can capture hub-oriented size preferences when the GCN encoder is initialized from degree ordering, the GCN's aggregation scheme is not expressive enough to adapt the full sequence to the global connectivity requirements of BA graphs.

\section{Attention RL Ordering on Barab\'asi--Albert Graphs}
\label{sec:results-attention-rl}

The final learned-ordering method was evaluated on Barab\'asi--Albert graphs using the same metrics as the deterministic baselines. Table~\ref{tab:ba-rl-final} reports the final generated-sample metrics for Attention RL ordering.

\begin{table}[H]
\centering
\footnotesize
\begin{tabular}{lrrrrr}
\toprule
\textbf{Method} & \textbf{Avg. size} & \textbf{Valid size} & \textbf{Connected} & \textbf{Degree Gini} & \textbf{Hubiness} \\
\midrule
Attention RL ordering & 41.4140 & 1.0000 & 0.9960 & 0.3081 & 0.3855 \\
\bottomrule
\end{tabular}
\caption{Final Attention RL ordering results on freshly sampled Barab\'asi--Albert graphs.}
\label{tab:ba-rl-final}
\end{table}

The most important result is that Attention RL ordering achieves a valid-size ratio of $1.0000$ and a connected ratio of $0.9960$. This places the learned method very close to the best connectivity behavior observed among the deterministic baselines, while also preserving perfect size validity. The average generated size is $41.414$, which is much closer to the target BA range than the smaller graphs produced by BFS. This indicates that the learned policy does not simply copy BFS connectivity behavior; it finds a different region of the trade-off space.

The degree Gini coefficient of $0.3081$ is lower than the strongest deterministic hub-oriented baselines, especially degree ordering and DFS. However, hubiness is $0.3855$, which is competitive with the strongest hubiness values and higher than degree, spectral, and LexBFS under the deterministic protocol. This suggests that the learned method preserves a dominant-hub signal even when the global degree inequality is less extreme. In practical terms, the generated graphs are connected and size-valid while still retaining a meaningful hub structure.

Table~\ref{tab:ba-comparison-full} provides a comprehensive comparison including the two learned-ordering configurations alongside the full set of deterministic baselines. The table makes explicit what a single deterministic ordering cannot achieve: no fixed heuristic simultaneously occupies the top positions in valid-size ratio, connectivity, and hubiness. The comparison with Degree ordering is particularly revealing: both methods achieve perfect valid-size ratio, but Attention RL ordering raises connectivity from $0.6490$ to $0.9960$ while also improving hubiness from $0.3326$ to $0.3855$. This shows that the learned policy does not merely reproduce the Degree initialization; reinforcement learning refines that hub-oriented prior into a substantially stronger global serialization strategy.

\begin{table}[H]
\centering
\footnotesize
\begin{tabular}{lrrrrr}
\toprule
\textbf{Method} & \textbf{Avg. size} & \textbf{Valid size} & \textbf{Connected} & \textbf{Degree Gini} & \textbf{Hubiness} \\
\midrule
BFS                       & 27.3110 & 0.8130 & 1.0000 & 0.2923 & 0.4282 \\
Degeneracy                & 37.0770 & 0.9670 & 0.9740 & 0.3287 & 0.3993 \\
MCS                       & 40.5150 & 0.9670 & 0.9110 & 0.3323 & 0.3886 \\
DFS                       & 40.8660 & 0.9630 & 0.8280 & 0.3411 & 0.4202 \\
Random                    & 46.2100 & 0.9940 & 0.7930 & 0.3328 & 0.3460 \\
Degree                    & 50.3650 & 1.0000 & 0.6490 & 0.3516 & 0.3326 \\
Spectral                  & 48.4210 & 0.9910 & 0.6010 & 0.3184 & 0.3085 \\
LexBFS                    & 50.2580 & 0.9540 & 0.4910 & 0.3358 & 0.3237 \\
\midrule
RL GCN ordering           & 39.3220 & 1.0000 & 0.6334 & 0.3272 & 0.3443 \\
\textbf{Attention RL ordering} & \textbf{41.4140} & \textbf{1.0000} & \textbf{0.9960} & \textbf{0.3081} & \textbf{0.3855} \\
\bottomrule
\end{tabular}
\caption{Comparison of BA results under deterministic node orderings and the two RL-learned orderings.}
\label{tab:ba-comparison-full}
\end{table}

The progression from RL GCN ordering to Attention RL ordering captures the impact of the three architectural changes introduced in the final configuration. The attention-based encoder brings transformer-style neighborhood weighting into the ordering decisions, allowing the policy to assign different importance to different neighbors depending on the current structural context. The deeper three-layer message-passing stack enlarges the receptive field available before each selection step, enabling decisions to reflect a broader portion of the graph topology. The degree-based warm start gives the controller an explicitly hub-aware prior from the beginning, after which reinforcement learning refines that prior toward a better global compromise between size, connectivity, and BA-style structural organization. Together these changes move the learned policy from a regime in which connectivity was the main weakness to one in which connectivity is nearly perfect.

\section{Loss Evolution Across Outer Iterations}
\label{sec:results-loss-evolution}

The validation and test loss curves in Figure~\ref{fig:ba-nll-evolution-loops} are essential for interpreting the learned method. They show the behavior of the coupled policy-generator system across outer iterations. Rather than training a generator once under a fixed sequence rule, the method repeatedly updates the ordering policy and retrains the generator on the resulting sequences. The decreasing loss trajectory indicates that the learned orderings become increasingly compatible with the downstream generator.

\begin{figure}[H]
    \centering
    \maybeincludegraphics[width=0.75\linewidth]{figures/loss_output_new_run.png}
    \caption{Validation and test loss across outer iterations of the Attention RL ordering procedure. The best validation checkpoint is reached at outer iteration~61, indicated by the dashed line. The trajectory remains oscillatory, as expected in an alternating RL--DGMG setting, but the late emergence of the best checkpoint shows that the policy continues to explore and refine the ordering space rather than collapsing prematurely.}
    \label{fig:ba-nll-evolution-loops}
\end{figure}

The Attention RL ordering remains competitive over a much longer optimization horizon than the GCN baseline and reaches its best validation checkpoint only at outer iteration 61. Although the trajectory remains oscillatory, as expected in an alternating RL--DGMG setting, the late best checkpoint indicates that the stronger policy continues refining the ordering space instead of collapsing prematurely to a suboptimal solution. The figure is shown up to outer iteration 62 rather than stopping exactly at the best point: iteration 61 is the selected checkpoint, while iteration 62 is included to make explicit the stopping logic based on the difference between consecutive loss values.

This loss behavior is important because structural metrics alone could be misleading. A generator might produce connected graphs by collapsing to a narrow subset of structures, or it might obtain good size control without learning the full action distribution. By monitoring validation and test loss, the experiment checks that the method also improves predictive fit. The final checkpoint is therefore supported by both likelihood-based and sample-based evidence.

\section{Interpretation}
\label{sec:results-interpretation}

The final results support the central claim of the dissertation in three ways. First, deterministic orderings strongly affect performance even when the generator architecture is unchanged. This establishes ordering as a real modeling variable. Second, the best deterministic ordering depends on topology and metric. This shows that no fixed rule should be assumed optimal in general. Third, Attention RL ordering learns a strong compromise on BA graphs, combining size validity, near-perfect connectivity, and meaningful hub structure.

The learned policy should be understood as a structural controller. It does not generate graphs directly; DGMG remains the graph generator. The policy controls how each training graph is exposed to the generator. By changing this exposure, it changes the conditional histories used for training and therefore changes what the generator learns. This separation is conceptually important because it means learned ordering can, in principle, be attached to other sequential graph generators that depend on graph-to-sequence conversion.

The results also show why the BA setting is harder than the earlier synthetic families. In cycles, trees, binary trees, and stars, one can define a precise validity predicate. In BA graphs, several structural properties matter simultaneously. A method that optimizes only one may harm another. Attention RL ordering is promising because it moves toward a balanced region of the metric space rather than maximizing a single structural statistic. Attention RL ordering is the only method in the comparison that simultaneously achieves perfect valid-size ratio, near-perfect connectivity, realistic graph size, and competitive hub-oriented structure, making it a genuinely strong alternative to handcrafted orderings.

\section{Limitations}
\label{sec:results-limitations}

The results should be interpreted with limitations in mind. The experiments use synthetic graph families, which are valuable for controlled analysis but do not capture the full complexity of molecular, social, or biological graph distributions. The learned-ordering method is also more computationally expensive than deterministic preprocessing because it requires an outer loop that couples policy learning with generator training. As noted in the experimental protocol, the full Attention RL run on BA graphs required approximately 7 hours on the MacBook M1 hardware used in this dissertation. In addition, the reward is based on validation and sample metrics, which means that its quality depends on the evaluation budget and on the choice of metric weights.

Another limitation is that the policy-gradient estimator can have high variance. The final method mitigates this issue with an attention-based architecture, additional graph layers, and a degree warm start, but further variance-reduction techniques such as actor-critic methods or importance sampling could improve stability. Finally, the learned policy is evaluated primarily on BA graphs in the final version of the dissertation. Extending the learned-ordering experiments to larger real-world datasets remains an important direction for future work.

\section{Chapter Summary}
\label{sec:results-summary}

This chapter showed that deterministic orderings produce large and interpretable differences in DGMG performance. The best ordering is topology-dependent: LexBFS and DFS are strongest for cycles, BFS is strongest for trees and binary trees, degree ordering is strongest for stars, and BA graphs exhibit trade-offs among connectivity, size, and hub structure. The GCN RL ordering confirmed the viability of the learned-ordering framework but did not overcome the connectivity limitation. The final Attention RL ordering method improves this trade-off on BA graphs, achieving perfect valid-size ratio, near-perfect connectivity, realistic average size, and competitive hubiness. These results support the dissertation's main thesis that node ordering should be treated as an optimization target in autoregressive graph generation.

%========================
% chapters/09_conclusion.tex
%========================
\chapter{Conclusion}
\label{chap:conclusion}

This dissertation studied node ordering as a central component of autoregressive graph generation. The starting point was a simple observation: sequential graph generators do not observe unordered graphs directly. They observe construction sequences, and these sequences depend on how the nodes are ordered. Therefore, an ordering is not just a preprocessing convention; it defines the conditional histories, action candidates, and structural regularities available to the generator during training.

The first contribution was a controlled analysis of deterministic orderings under a DGMG-style generator. By keeping the generator fixed and changing only the graph-to-sequence map, the dissertation isolated the effect of ordering. The results showed that orderings induce strong and topology-dependent differences. Cycles favor path-like orderings such as DFS and LexBFS. Trees and binary trees favor BFS because it preserves parent-frontier locality. Stars favor degree ordering because early hub exposure simplifies most destination decisions. Barab\'asi--Albert graphs reveal a more complex trade-off, in which different orderings emphasize size, connectivity, degree inequality, or hub dominance. These findings support the interpretation that ordering acts as an entropy-shaping mechanism for the sequential decision problem.

The second and main contribution was a learned-ordering framework based on reinforcement learning. The dissertation formulated node sequencing as a graph-conditioned Markov Decision Process in which a policy selects one unselected node at a time until a complete ordering is produced. The resulting ordering is translated into a DGMG action sequence, and the reward is computed from downstream generator behavior. The reward design explicitly combines predictive fit, measured as negative validation log-likelihood, with structural validity metrics computed on freshly sampled generated graphs. Because the ordering decisions are discrete and their effect is delayed, the policy is optimized with a policy-gradient estimator. The theorem presented in Chapter~\ref{chap:rl-ordering} formalizes this gradient and justifies the update used by the method.

An intermediate GCN-based RL policy confirmed that the general learned-ordering framework is viable and capable of achieving strong size control on BA graphs. However, connectivity remained a weakness, motivating the architectural upgrades introduced in the final configuration. The final version of the method, Attention RL ordering, strengthened the policy with an attention-based graph encoder, three graph layers, and a degree-based warm start. The method was evaluated on Barab\'asi--Albert graphs, where it achieved perfect valid-size ratio, near-perfect connectivity, realistic average size, and competitive hubiness. The result is significant because it combines properties that are separated across deterministic baselines: BFS is strong for connectivity but weak for size, while degree ordering is strong for size and hub exposure but weak for connectivity. The learned method moves toward a hybrid region of the metric space, which is precisely the motivation for replacing a fixed heuristic with an adaptive policy.

The experiments are based on synthetic graph families, and although these families are useful for controlled interpretation, real-world graph distributions may introduce additional constraints such as labels, attributes, edge types, and domain-specific validity. The learned-ordering method also has higher computational cost than deterministic ordering and depends on reward design. Policy-gradient optimization can be noisy, especially when the reward is based on generator training and sample evaluation. These limitations do not undermine the main conclusion, but they define the scope in which the results should be generalized.

Several directions remain open. The first is to evaluate learned ordering on real datasets such as molecular graphs, program graphs, or biological networks. The second is to combine learned ordering with stronger autoregressive backbones, including transformer-based graph generators such as GraphGPT \cite{zhao2023graphgpt} and G2PT \cite{chen2025g2pt}. The third is to improve the reinforcement-learning component with variance reduction, actor-critic methods, imitation from strong heuristics, or curriculum learning over graph size. The fourth is to study whether the learned policy discovers interpretable hybrid strategies, for example by measuring agreement with BFS, degree ordering, or MCS at different stages of the trajectory. The policy-entropy diagnostic shown in Figure~\ref{fig:rl-entropy-mean} suggests that the policy becomes increasingly concentrated over training, which is consistent with a transition from broad exploration to confident ordering decisions; characterizing when and why this concentration occurs could yield useful insights for improving training stability.

Overall, the dissertation argues that the ordering problem should be treated as part of the model design in sequential graph generation. A graph generator trained on ordered sequences is only as effective as the representation through which the graph is exposed. Handcrafted orderings provide useful priors, but they are topology-dependent and metric-dependent. Learned orderings offer a way to adapt this prior to the data and to the downstream generator. This perspective opens a path toward graph generators in which the construction order is not assumed in advance, but learned as part of the generative process itself.


\bibliographystyle{ieeetr}
\bibliography{refs}

\end{document}
